\Needspace{10\baselineskip}
\section{Relative projection and supersolutions}
\label{sec:big-relative-projection}

To pass from the minimal class to the reference $\varphi$, we project a solution $z\in\mathcal E_V$ onto $\mathcal E_\varphi$. The projection is the increasing limit $P_\theta[\varphi](z)$ of the rooftops $P_\theta(\varphi+C,z)$ as $C\to\infty$, in the notation of \cref{lem:relative-obstacle-contact}. The small stretches of \cref{lem:positive-volume-stretch} first show that the rooftop $P_\theta(\varphi,z)$ is not identically $-\infty$ and has the full upper volume of $\varphi$. Contact support then bounds the measure of the projection by the measure of $z$ on the set where the two coincide. The argument uses only that the reference of $z$ lies above $\varphi$, so we formulate it for an arbitrary ordinary model above $\varphi$; this will also allow the reference to vary.

Let $\eta\in\PSH(X,\theta)$ satisfy $P_\theta[\eta]=\eta\geq\varphi$, and write $\mathcal E_\eta=\mathcal E_{\mathrm{env}}(X,\theta,\eta)$.

\begin{proposition}\label{prop:big-cross-class-roof}
For every $z\in\mathcal E_\eta$, the rooftop $P_\theta(\varphi,z)$ is not identically $-\infty$ and belongs to $\mathcal E_\varphi$.
\end{proposition}
\Needspace{7\baselineskip}
\begin{proof}
First suppose $z\leq\eta$; recall also that $\varphi\leq\eta\leq0$. For $b>1$ sufficiently close to one, \cref{lem:positive-volume-stretch} gives
\[
 h_b=Q_b(\varphi,\eta)\leq\varphi,
\]
which is not identically $-\infty$. Set $a=1-b^{-1}$. The stretching lemma \cref{lem:relative-stretch-closure}, now applied in $\mathcal E_\eta$, gives $z_{1/a}=Q_{1/a}(z,\eta)\in\mathcal E_\eta$. The potential
\[
 k_a=a z_{1/a}+(1-a)h_b
\]
lies below both $\varphi$ and $z$. Indeed, $z_{1/a}\leq\eta$ and $(1-a)h_b+a\eta\leq\varphi$ give the first bound, while $h_b\leq\eta$ and $a z_{1/a}+(1-a)\eta\leq z$ give the second. Thus $p=P_\theta(\varphi,z)$ is not identically $-\infty$. Volume addition and monotonicity give
\[
 (1-a)^n\overline{\operatorname{vol}}_\theta(h_b)
 \leq\overline{\operatorname{vol}}_\theta(k_a)
 \leq\overline{\operatorname{vol}}_\theta(p)
 \leq\overline{\operatorname{vol}}_\theta(\varphi).
\]
As $b\downarrow1$, we have $a\downarrow0$, and the first term tends to the last by \cref{lem:positive-volume-stretch}. The upper volumes of $p$ and $\varphi$ are therefore equal. Since $p\leq\varphi$, \cref{thm:relative-equivalence} gives $p\in\mathcal E_\varphi$.

\Needspace{6\baselineskip}
For general $z$, the envelope identity in $\mathcal E_\eta$ gives $z-\sup_Xz\leq\eta$. Set $S=\max(\sup_Xz,0)$. The preceding case gives $P_\theta(\varphi,z-S)\in\mathcal E_\varphi$, and
\[
 P_\theta(\varphi,z-S)\leq P_\theta(\varphi,z)\leq\varphi.
\]
Upward stability finishes the proof.
\end{proof}

\Needspace{12\baselineskip}
\begin{proposition}\label{prop:big-relative-projection}\index{Envelope!relative projection}
Let $z\in\mathcal E_\eta$ and let $Q=P_\theta[\varphi](z)=\bigl(\lim_{C\to\infty}P_\theta(\varphi+C,z)\bigr)^*$ be its projection onto the relative class of $\varphi$.
Then $Q\in\mathcal E_\varphi$, $Q\leq z$, and
\[
 \mu_Q\leq\mathbf1_{\{Q=z\}}\mu_z.
\]
\end{proposition}
\begin{proof}
For $C\geq0$ put $p_C=P_\theta(\varphi+C,z)$. The preceding proposition gives $p_0\in\mathcal E_\varphi$. Since $p_0\leq p_C\preceq\varphi$, upward stability gives $p_C\in\mathcal E_\varphi$. The envelope identity, together with $p_C\leq z$, yields the uniform bound
\[
 p_C\leq\varphi+\sup_Xp_C\leq\varphi+\sup_Xz.
\]
Consequently the increasing regularized limit $Q$ is not identically $-\infty$, belongs to $\mathcal E_\varphi$ and is at most $z$. The monotone convergence holds in capacity by \cref{lem:monotone-psh-capacity}.

Finally, \cref{lem:relative-obstacle-contact}, applied with $u=\varphi$ and $w=z$, gives
\begin{equation}\label{eq:big-projection-contact-tail}
 \mu_Q\leq\mathbf1_{\{Q=z\}}\mu_z.
\end{equation}
\end{proof}

\Needspace{16\baselineskip}
\begin{corollary}\label{cor:big-relative-supersolution}
For every nonzero finite non-pluripolar measure $\nu$, there is $Q\in\mathcal E_\varphi$ such that
\[
 0\ne \mathrm{e}^{-Q}\mu_Q\leq\nu.
\]
\end{corollary}
\begin{proof}
By \cref{thm:big-minimal-general-measures}, there is $z\in\mathcal E_V$ with $\mu_z=\mathrm{e}^z\nu$. Since $V$ is an ordinary model and $V\geq\varphi$, the preceding proposition applies with $\eta=V$ and gives $Q=P_\theta[\varphi](z)\in\mathcal E_\varphi$. On the contact set $\{Q=z\}$,
\[
 \mu_Q\leq\mathbf1_{\{Q=z\}}\mathrm{e}^z\nu\leq \mathrm{e}^Q\nu.
\]
The weights are finite off the pole sets, which carry no mass. Thus $\mathrm{e}^{-Q}\mu_Q\leq\nu$, and this measure is nonzero because $\mu_Q(X)\geq\ell>0$.
\end{proof}

\subsection{Varying model references}
\label{subsec:big-varying-models}

We now let the reference vary. Recall from \cref{def:model-volume-comparison} the comparison
\[
 d_\theta(u,v)=2\,\overline{\operatorname{vol}}_\theta(\max(u,v))
 -\overline{\operatorname{vol}}_\theta(u)-\overline{\operatorname{vol}}_\theta(v)
\]
and, for $\delta>0$, the set $\mathcal M_\delta(\theta)$ of ordinary models $u=P_\theta[u]$ with $\underline{\operatorname{vol}}_\theta(u)\geq\delta$. For models $\varphi_j\to\varphi$ in this comparison, \cref{thm:model-volume-capacity-convergence} provides increasing lower models whose upper volumes approach that of $\varphi$. Projecting finite families of potentials onto these lower models makes their rooftops comparable. To keep the limit in the relative class of $\varphi$, domination of the measures is needed only where the potentials lie below one fixed member of that class; domination on all of $X$ gives, in addition, convergence in capacity and convergence of the measures.

\begin{theorem}\label{thm:big-varying-model-closure}\index{Relative full-mass class!varying models}
Let $\varphi_j,\varphi\in\mathcal M_\delta(\theta)$ for some $\delta>0$, with $d_\theta(\varphi_j,\varphi)\to0$. Suppose $u_j\in\mathcal E_{\varphi_j}$, $u_j\to u\in\PSH(X,\theta)$ in $L^1(X)$, and $\chi\in\mathcal E_\varphi$. If one finite non-pluripolar measure $\nu$ satisfies
\[
 \mathbf1_{\{u_j\leq\chi\}}\mu_{u_j}\leq\nu,
\]
then $u\in\mathcal E_\varphi$. Moreover, after extraction and relabeling, the potentials
\[
 \psi_j=P_\theta\left(\inf_{l\geq j}\varphi_l\right),\qquad
 h_j=P_\theta\left(\inf_{l\geq j}u_l\right)
\]
satisfy $\psi_j\in\mathcal M_{\delta/2}(\theta)$ and $h_j\in\mathcal E_{\psi_j}$.
If $\mu_{u_j}\leq\nu$ on $X$, the original sequence converges to $u$ in capacity and $\mu_{u_j}\rightharpoonup\mu_u$.
\end{theorem}
\begin{proof}
We first construct the lower models $\psi_j$, then a common lower bound for rooftops of the $u_l$ in the classes $\mathcal E_{\psi_j}$, and finally treat domination on all of $X$.

By Hartogs' lemma, $u_j$, $u$ and $\chi$ have a common upper bound; translate them to be nonpositive. The envelope identity then gives $u_j\leq\varphi_j$. Since $\varphi_j\to\varphi$ in $L^1$ by \cref{prop:model-volume-separation}, it follows that $u\leq\varphi$. Put $\Lambda=\nu+\mu_\chi$. Using \cref{lem:fixed-npp-integrals}, extract a subsequence along which $u_j\to u$ almost everywhere for smooth volume and
\[
 I_j:=\int_X|\mathrm{e}^{u_j}-\mathrm{e}^u|\,\mathrm{d}\Lambda\leq2^{-j}.
\]
Apply \cref{thm:model-volume-capacity-convergence} to extract further, keeping these estimates, and obtain increasing models $w_j\leq\varphi_j,\varphi$ with lower volume at least $\delta/2$ and $d_\theta(w_j,\varphi)\to0$. For $l\geq j$ we have $w_j\leq w_l\leq\varphi_l$. Hence $\psi_j$ is not identically $-\infty$, and
\[
 w_j\leq\psi_j\leq\varphi.
\]
The upper bound follows by passing to the $L^1$ limit in $\psi_j\leq\varphi_l$ as $l\to\infty$. Monotonicity of the ordinary envelope gives $P_\theta[\psi_j]\leq\varphi_l$ for all $l\geq j$, so $\psi_j$ is an ordinary model. It increases with $j$, has lower volume at least $\delta/2$, and
\begin{equation}\label{eq:varying-reference-volume-limit}
 \overline{\operatorname{vol}}_\theta(\psi_j)
 \longrightarrow\overline{\operatorname{vol}}_\theta(\varphi).
\end{equation}

Set $g_j=P_\theta[\psi_j](\chi)$. \Cref{prop:big-relative-projection}, with $(\varphi,\eta)$ replaced by $(\psi_j,\varphi)$, gives $g_j\in\mathcal E_{\psi_j}$, $g_j\leq\chi$ and $\mu_{g_j}\leq\mu_\chi$. Moreover $g_j$ increases with $j$. For $N\geq j$, consider
\[
 q_{j,N}=P_\theta(g_j,u_j,\ldots,u_N).
\]
For each $l\geq j$, the projection $P_\theta[\psi_j](u_l)$ belongs to $\mathcal E_{\psi_j}$ by the same proposition, now with upper reference $\varphi_l$. The finite rooftop of these projections and $g_j$ lies below $q_{j,N}$. Since $q_{j,N}\leq g_j$, the finite-rooftop proposition and upward stability give $q_{j,N}\in\mathcal E_{\psi_j}$, and hence $\mu_{q_{j,N}}(X)\geq\delta/2$.

Partition the contact set of $q_{j,N}$ into disjoint pieces, one for each obstacle. On a piece where $q_{j,N}=u_l$, we have $u_l\leq g_j\leq\chi$, so the hypothesis bounds the measure there by $\nu$. On the piece where $q_{j,N}=g_j$, the measure is bounded by $\mu_{g_j}\leq\mu_\chi$. Consequently $\mu_{q_{j,N}}\leq\Lambda$. Put $F_j=\min(\mathrm{e}^u,\mathrm{e}^{g_j})$. On the $u_l$ piece,
\[
 |\mathrm{e}^{q_{j,N}}-F_j|\leq|\mathrm{e}^{u_l}-\mathrm{e}^u|.
\]
On the $g_j$ piece, $g_j=q_{j,N}\leq u_N$, so the error is
$(\mathrm{e}^{g_j}-\mathrm{e}^u)_+\leq|\mathrm{e}^{u_N}-\mathrm{e}^u|$.
The $u_N$ piece and the $g_j$ piece are controlled by the two summands $\nu$ and $\mu_\chi$ of $\Lambda$, so together they contribute at most $I_N$. Thus
\begin{equation}\label{eq:varying-roof-exponential-error}
 \int_X|\mathrm{e}^{q_{j,N}}-F_j|\,\mu_{q_{j,N}}
 \leq\sum_{l=j}^N I_l\leq2^{1-j}.
\end{equation}

Since $F_j\geq F_1>0$ outside a pluripolar set, we may choose $\epsilon>0$ with $\Lambda(\{F_1\leq\epsilon\})<\delta/4$. The lower mass bound and the domination by $\Lambda$ give
\[
 \epsilon\delta/4
 \leq\int_X F_j\,\mu_{q_{j,N}}
 \leq2^{1-j}+\Lambda(X)\mathrm{e}^{\sup_Xq_{j,N}}.
\]
For every sufficiently large $j$, this bounds $\sup_Xq_{j,N}$ from below, uniformly in $N$. Hence the decreasing limit $q_j$ as $N\to\infty$ is not identically $-\infty$. Applying \cref{thm:big-relative-fixed-continuity} in the fixed class $\mathcal E_{\psi_j}$, with dominating measure $\Lambda$, gives $q_j\in\mathcal E_{\psi_j}$. Now
\[
 q_j\leq h_j\leq\psi_j,
\]
so $h_j\in\mathcal E_{\psi_j}$ by upward stability. We discard the finitely many earlier indices. Since $h_j\leq u\leq\varphi$, monotonicity and \eqref{eq:varying-reference-volume-limit} show that $u$ and $\varphi$ have the same upper volume, and \cref{thm:relative-equivalence} gives $u\in\mathcal E_\varphi$.

Suppose now that $\mu_{u_j}\leq\nu$ on $X$. The finite rooftops $p_{j,N}=P_\theta(u_j,\ldots,u_N)$ are not identically $-\infty$, since they lie above $h_j$. Partitioning their contact sets gives
\[
 \mu_{p_{j,N}}\leq\nu,\qquad
 \int_X|\mathrm{e}^{p_{j,N}}-\mathrm{e}^u|\,\mu_{p_{j,N}}\leq2^{1-j}.
\]
For fixed $j$, they decrease to $h_j$ as $N\to\infty$. These lower envelopes increase to a regularized limit $h\leq u\leq\varphi$, and their upper volumes tend to that of $\varphi$; hence $h\in\mathcal E_\varphi$. Both monotone limits hold in capacity by \cref{lem:monotone-psh-capacity}. Applying \cref{lem:weighted-extended-lsc} twice gives
\[
 \int_X|\mathrm{e}^h-\mathrm{e}^u|\,\mu_h=0.
\]
Hence $\mu_h(\{h<u\})=0$, and zero domination gives $h=u$. The upper envelopes $(\sup_{k\geq j}u_k)^*$ decrease to $u$ by Hartogs' lemma and the almost-everywhere convergence, and squeezing gives convergence in capacity. Finally,
\[
 \mu_{u_j}(\{u_j\leq-T\})\leq\nu(\{h_1\leq-T\})\longrightarrow0
\]
uniformly in $j$, and the tail criterion \cref{thm:general-background-damping} gives weak convergence of the measures. Since the extraction can start from any subsequence, both conclusions hold for the original sequence.
\end{proof}
